\documentclass{article}
\usepackage{icmlko}
\usepackage{graphicx}
\usepackage{booktabs}
\usepackage{amsmath}

\title{같은 하늘 아래로 바닥을 옮겼더니 신호가 그 위에 있었다 \\
  {\large 그런데 그 빈틈의 절반은 관측창의 가장자리였고, 나머지의 대부분은 달력이었다}}

\begin{document}
\maketitle

\begin{abstract}
직전 노트는 격자 위에서 사건 간격을 재어 개선 $+0.7545$ 일을 얻고도,
옆에 세운 순열 바닥이 $+1.0693$ 으로 더 높아 \textbf{③ 은 0} 이라고 적었다.
프레시 티처가 그 바닥을 반증했다 --- \textbf{순열이 사건을 옮기면서 급증까지 지웠고},
두 세계의 기후값 MAE 가 \textbf{19.319\%} 달랐다.
이 노트는 바닥을 \textbf{같은 세계 안으로} 옮긴다.
\claimx{\textbf{간격순서 순열}(격자 안에서 간격 수열만 뒤섞기)은 진짜 세계와 아홉 가지를 공유한다.}
{하나라도 어긋나면 이 귀무는 못 쓴다. 실측 --- 사건 185{,}070 $\cdot$ 간격 176{,}770 $\cdot$
홀드아웃 간격 54{,}862 $\cdot$ 군집 2{,}277 $\cdot$ 학습 121{,}908 이 \textbf{전부 같고},
mag 3 분위와 간격 중앙값$\cdot$평균이 \textbf{소수점까지 같으며},
격자별 간격 다중집합이 \textbf{7{,}537 격자 전수에서 하나도 안 깨졌다}.
기후값 MAE 상대차는 \textbf{0.4727\%} 로 문턱 5\% 안이다.}
\claimx{그 바닥 위에서 \textbf{신호가 바닥보다 높다} --- $+0.7545$ 대 $+0.2827$, 두 구간이 안 겹친다.}
{겹치면 이 문장은 무너진다. 빈틈이 \textbf{0.3122 일}이다.}
🔴 \textbf{바뀐 것은 자가 아니라 바닥이 놓인 세계다.}
직전 노트의 진짜 비교값은 \textbf{소수점 끝까지 그대로}이고, 옛 귀무도 그대로 재현된다.
\claimx{🔴 \textbf{그리고 프레시 티처가 그 빈틈의 이름을 무너뜨렸다.}}
{빈틈이 「순서를 부수면 사라지는 것」이라면 달력과 관측창 가장자리를 통제해도 남아야 한다.
실측 --- 달력을 두 세계에서 없애고 격자별 \emph{첫} 간격을 빼면
진짜 $0.2165\,[0.1980, 0.2368]$ 대 N2 $0.1986\,[0.1777, 0.2194]$ 로 \textbf{두 구간이 겹친다}.
빈틈의 \textbf{77.8\%가 달력 결합}이고, \textbf{50.9\%가 홀드아웃의 4.150\%에 해당하는
2{,}277 행}(격자마다 첫 간격 하나)에서 온다.}
그러므로 이 노트가 실제로 보인 것은 \textbf{③ 의 0.5 칸}이다 ---
\textbf{바닥이 신호 아래로 내려온 것은 참}이고, \textbf{그 빈틈을 「시간 구조」라 부른 것은 틀렸다}.
\end{abstract}

\section{무엇이 문제였나}

「바닥」이라는 말에는 조용한 전제가 있다 --- \textbf{바닥과 신호가 같은 세계에서 잰 수}라는 것.
직전 노트의 순열은 격자 안에서 사건 \emph{날짜}만 재추첨했고 급증 배수 \texttt{r} 은 원래 자리에 두고 왔다.
그래서 귀무의 「사건」은 대개 급증이 아니었고, 두 세계의 난이도가 벌어졌다.
\textbf{서로 다른 난이도의 두 수를 일 단위로 나란히 놓은 것}이 조항 60 이 금지하는 형태다.

\section{바닥을 셋 만들고 관문에 걸었다}

\begin{figure}[h]\centering
\includegraphics[width=\linewidth]{figs/fig1.pdf}
\caption{(a) 네 세계의 개선과 BCa --- 흐린 둘은 비교 불가라 판정에 안 썼다.
(b) 비교가능성 관문이 둘을 떨어뜨렸다. (c) N2 가 진짜와 공유하는 아홉 --- 다른 것은 \textbf{간격의 순서} 하나다.}
\end{figure}

\begin{table}[h]\centering\small
\begin{tabular}{lrrrl}
\toprule
세계 & 개선(일) & BCa 95\% & 개선율 & 비교가능성 \\
\midrule
\textbf{진짜} & $\mathbf{+0.7545}$ & $[+0.6436, +0.8861]$ & 6.221\% & --- \\
\textbf{N2 간격순서 순열} & $\mathbf{+0.2827}$ & $[+0.2375, +0.3314]$ & 2.320\% & ✅ \textbf{0.4727\%} \\
N1 「\texttt{r} 도 섞기」 & $+0.9799$ & $[+0.8868, +1.0879]$ & 6.771\% & ✗ 19.319\% \\
N0 옛 귀무 & $+1.0693$ & $[+0.9665, +1.1833]$ & 7.389\% & ✗ 19.319\% \\
\bottomrule
\end{tabular}
\end{table}

구간은 전부 \texttt{lab.pairboot.cluster\_boot} \textbf{BCa}(B$=$2000, 군집$=$격자)이고 \textbf{폴백 0 건}이다(규약 47).

\subsection{🔴 티처의 진단은 증상을 맞혔고 병은 다른 것이었다}

이슈의 처방 1 은 「\texttt{r} 도 섞어라」였다. 그대로 하니 mag 3 분위는 고쳐졌다
($[1.28950, 1.61911] \approx$ 진짜 $[1.29092, 1.62410]$).

\claimx{\textbf{그런데 난이도 차는 하나도 안 줄었다} --- 상대차가 \textbf{19.319\%} 로 옛 귀무와 소수점까지 같다.}
{급증 크기가 원인이었다면 이 수가 줄었어야 한다. 안 줄었다.
기후값 팔은 \texttt{mag} 을 아예 안 보기 때문이다 --- 난이도를 만든 것은 \textbf{간격 분포}였다
(귀무 중앙값 7.0$\cdot$평균 18.108 대 진짜 6.0$\cdot$16.009).}
그러므로 이 노트를 연 것은 처방 1 이 아니라 \textbf{처방 2(비교가능성 관문 W7)}다 ---
\textbf{고치라는 지시가 아니라, 고쳐졌는지 재는 자}가 일을 했다.

\section{판정}

\begin{figure}[h]\centering
\includegraphics[width=\linewidth]{figs/fig2.pdf}
\caption{(a) 같은 신호에 바닥만 바꾸면 결론이 뒤집힌다.
(b) 배선과 관문 --- 검정력 옆에 \textbf{분모}를 적었다.}
\end{figure}

\claimx{\textbf{같은 세계의 바닥이 신호 아래로 내려왔다} --- $+0.2827$ 대 $+0.7545$, 안 겹친다.}
{이 문장은 서고, 티처가 러너를 다시 돌려 \textbf{산출물 키 530 개 중 차이 0 개}로 재현했다.}

\subsection{🔴 정정 --- 그 빈틈의 이름이 틀렸다}

초판은 여기에 「나머지 $+0.4718$ 은 순서를 부수면 사라진다 --- 그것이 시간 구조의 몫이다」를 적고
\textbf{③ 이 0 을 벗어난다}로 헤드라인을 걸었다. \textbf{그 이름이 틀렸다.}

\begin{table}[h]\centering\small
\begin{tabular}{lrrl}
\toprule
절제 & 진짜 & N2 & 겹치나 \\
\midrule
원판 & $0.7545\,[0.6436, 0.8861]$ & $0.2827\,[0.2375, 0.3314]$ & ✗ \\
첫 간격 제외 & $0.4479\,[0.3757, 0.5272]$ & $0.2062\,[0.1818, 0.2322]$ & ✗ \\
달력 제거 & $0.3238\,[0.2837, 0.3674]$ & $0.2190\,[0.1941, 0.2433]$ & ✗ \\
\textbf{둘 다} & $\mathbf{0.2165\,[0.1980, 0.2368]}$ & $\mathbf{0.1986\,[0.1777, 0.2194]}$ & \textbf{✓ 겹친다} \\
\bottomrule
\end{tabular}
\end{table}

\claimx{빈틈의 \textbf{77.8\%가 달력 결합}이고, \textbf{50.9\%가 홀드아웃의 4.150\%}에서 온다.}
{그 4.150\% 는 \textbf{2{,}277 행}이고 이것은 \textbf{홀드아웃 격자 수와 정확히 같다} ---
격자마다 \emph{첫} 간격 하나다. 그 행의 평균 간격이 $58.54$ 일로 나머지 $13.73$ 일의 4.26 배이고,
그 행만의 개선이 $+7.836\,[6.288, 9.410]$ 이다.
\textbf{577 일 관측창의 가장자리}이지 과정의 시간 구조가 아니다.}

\claimx{그러므로 이 노트가 실제로 민 것은 \textbf{③ 의 0.5 칸}이다.}
{「바닥이 신호 아래로 내려왔다」는 서고, 「그 빈틈이 시간 구조다」는 안 선다.
철회가 아니라 \textbf{재명명}이다.}

직전 노트의 두 문장은 이렇게 고쳐진다:
① 「바닥이 신호보다 위다」 $\to$ \textbf{아니다. 같은 세계의 바닥은 신호의 37.5\% 자리다.}
② 「얻은 것은 발생률뿐이다」 $\to$ \textbf{반만 맞다} --- 다만 나머지의 정체는
\textbf{시간 구조가 아니라 달력과 관측창 가장자리}다.

\subsection{🔴 그리고 관문이 통과시킨 이유가 틀렸다}

비교가능성 관문(기후값 MAE 상대차 $\le$ 5\%)의 \textbf{도달 가능 최댓값이 0.4731\%}다 ---
진짜 세계의 기후값 팔이 무정보 팔보다 나은 정도가 그것뿐이기 때문이다.
문턱은 그 \textbf{10.6 배}다.

\claimx{그러므로 실측 0.4727\%는 「거의 안 변했다」가 아니라 \textbf{「기후값 팔이 무정보 팔로 완전히 퇴화했다」}는 뜻이다.}
{반례가 있다 --- \texttt{mag} 짝만 깨는 귀무는 상대차가 \textbf{정확히 0.0\%}로 만점 통과하는데
개선이 $+0.5737\,[0.4832, 0.6845]$ 로 진짜와 \textbf{겹친다}.}
\textbf{비교가능성 결론 자체는 산다} --- 다만 근거가 이 관문이 아니라
\textbf{홀드아웃 간격 다중집합이 비트로 동일}하다는 쪽이다. \textbf{옳은 결론을 틀린 관문으로 정당화했다.}

\subsection{자기시험 --- 관이 같다는 것을 먼저 보였다}

새 러너가 직전 노트의 정본을 \textbf{소수점 끝까지} 재현한다:
$+0.7545204330866538$ $\cdot$ BCa $[0.6436055429260724, 0.8861167731832372]$ $\cdot$
홀드아웃 간격 54{,}862 $\cdot$ 군집 2{,}277 $\cdot$ 사건 185{,}070 $\cdot$ 간격 176{,}770.
\textbf{옛 귀무도 $+1.0693284969559986$ 로 정확히 재현}됐다.
입력 \texttt{daily.npz} 와 직전 모듈의 sha256 이 비트 일치한다.

배선 W0 은 여덟 열(\texttt{gi}$\cdot$\texttt{tk}$\cdot$\texttt{gap}$\cdot$\texttt{dow}$\cdot$
\texttt{quarter}$\cdot$\texttt{mag}$\cdot$\texttt{prevmed}$\cdot$\texttt{lvl})이 전부 같다고 하고,
\texttt{mag} 을 한 칸 굴려 심으니 \textbf{붉어졌다}.

\subsection{관문이 이번엔 제 일을 한다}

직전 노트의 정지 관문은 발화해도 계산을 안 멈췄다(\texttt{break} 가 바깥 절을 못 막았다).
계수를 $10^{-9}$ 로 낮춰 강제 발화시키니 \textbf{멈췄고}, 산출물의
\textbf{진짜 세계 개선 절이 0 개}다(직전 노트는 1 개였다).
그리고 검정력을 \textbf{5/5 = 1.000} 이라 적으면서 \textbf{「분모는 5 이고 관문 전체가 아니다」}를 같이 적었다.

\section{한계 --- 「없다」가 아니라 「못 했다」}

\begin{itemize}
\item 🔴 \claimx{\textbf{「N2 가 시각 구조 \emph{만} 부수는가」를 못 심어서 검정력 0 으로 신고했다.}}
      {간격 순서를 뒤섞으면 \textbf{달력과의 짝도 같이 깨진다}.
      그러므로 $+0.4718$ 을 「시간 구조의 몫」이라 부르는 것은 \textbf{상한}이지 분해가 아니다.}
\item \textbf{채택이 아니라 후보다}(노트 133). 첫 양수이고 프레시 티처를 아직 안 거쳤다
\item $\tau$ 민감도(1.10$\cdot$1.50)와 「큰 격자만」 절을 \textbf{안 돌렸다}
\item \textbf{귀무 세계의 덮음을 안 냈다.} 진짜 세계 덮음은 명목 0.80 대 실측 0.9047/0.8796 로
      \textbf{구간이 눈금을 안 맞춘다} --- 점추정은 유효하되 구간 예측은 아직 못 판다
\item ③ 의 \textbf{「어떤」은 안 밀었다} --- 이 자료에 사건의 \emph{사유}가 없다.
      민 것은 「언제」뿐이다
\item 직전 노트의 마스크가 버린 칸(관측된 (격자,날짜) 4{,}926{,}445 대 4{,}803{,}854)은
      이 노트도 \textbf{안 되살렸다}
\item 🔴 \textbf{③ 의 자는 둘인데 하나만 넘었다.} 「간격 예측이 기후값을 이기나」와
      \textbf{「구간 덮음」}이 둘 다 자인데, 덮음은 명목 0.80 대 실측 0.9047/0.8796 로
      \textbf{둘 다 $\pm$0.05 밖}이다. 커밋 제목은 넘은 하나만 말했다
\item 🔴 \textbf{바닥이 0 이 아니다.} 사전등록 §6-1(「귀무가 0 을 문다」)은 발화한 적 없고
      밟힌 것은 §6-3 이다. 시간 구조가 전혀 없는 세계에서도 이 자는 $+0.2827$ 로 유의하게 이긴다 ---
      기후값 팔이 \textbf{격자를 모르는 팔}이라서 발생률만으로 이미 이긴다(그 이득의 92.8\%가 \texttt{prevmed})
\item \textbf{N2 가 정말 섞였나를 심는 검사가 없다} --- 86.15\%는 자기신고이고,
      항등 N2 를 손으로 넣어도 다중집합 검사는 「통과」를 낸다
\item 러너가 \textbf{완주 두 번 뒤에 마지막 수정}됐다(러너 mtime 이 산출물보다 37 초 앞).
      🔴 다만 \textbf{문턱 상수는 모든 실행보다 앞}서 있어 결과를 보고 문턱을 고치지는 않았다
\end{itemize}

\end{document}
